\startreportpage{25.4}
\templaterow{结\quad 论}{70}{
四题均达到100\%完成度，题1至题4分别通过9/9、18/18、24/24、20/20项配置检查。题4对照评分核对并纠正等容量LAN子网与接口的对应关系后，全部评分项均为Correct。

题1表明，本地通信与默认网关配置可以分别验证：网关留空或填入不存在的本网段地址时，本地仍4/4成功，远端均0/4；恢复正确网关后，远端复测4/4成功。跨网段网关补测仍可达，且两个网关地址映射为同一MAC，现象与代理ARP机制一致，但本题未直接核对该功能开关。

题2的三台Fa0/0接口均为up/up，三台PC到各自网关均4/4成功。题3、题4均完成九个子网的划分及十二个接口的编址核对；这些结果证明配置符合活动检查要求，不据此推断各LAN间已实现端到端互通。
}
\templaterow{小\quad 结}{109.59}{
本次按“地址规划—逐项配置—读取实际状态—测试并保留结果”的顺序完成操作。IPv4地址、掩码和默认网关是相互关联的配置项，应先按位与求出网络地址，再判断目的地址是否需要交给默认网关，不能只凭地址前几位相似作判断。

网关对照测试中，留空与填写192.168.1.254都表现为远端超时，只有结合Default Gateway截图才能区分原因。正确网关组首次远端Ping为3/4，报告保留首包超时，并将恢复后的4/4复测作为另一条记录。跨网段网关仍可达的现象也提醒我：单次Ping成功不能代替地址规划和配置检查，应继续结合ARP表与活动评分核对。

VLSM划分需要同时满足容量与边界对齐。先按主机需求选择满足$2^h-2\geq H$的最小主机位数，再从大块到小块分配；串口两端还要核对属于同一/30，且不使用网络地址或广播地址。统计已分配空间时必须计入整个地址块，不能只累加可用主机数。

题4首次检查为14/20。对照评分明细发现，六个LAN中同容量子网与接口的对应关系反向；按活动接口映射纠正后达到20/20。这说明容量满足、子网不重叠还不够，必须逐项核对用途、子网和接口的一一对应关系。

配置检查、接口状态与Ping测试各有作用，未执行的跨LAN连通测试不写成成功。原始截图、命令记录及裁剪来源在本地保存，以便复核。
}
\TemplateReview
